\section{System Architectural Models (Modelli OOA)}

La presente sezione descrive l'architettura dinamica e strutturale del sistema \textbf{MyAma} mediante i modelli di \textit{Object Oriented Analysis} (OOA) sviluppati con il tool Visual Paradigm.

% =========================================================================
% 5.1 ACTIVITY DIAGRAMS
% =========================================================================
\subsection{Activity Diagrams}
Gli Activity Diagram descrivono i flussi operativi, le diramazioni decisionali e i punti di sincronizzazione nei processi cardine della piattaforma.

\subsubsection{Activity Diagram: Prenotazione Ritiro a Domicilio (Cittadino)}
Il diagramma modella il processo di inserimento dati rifiuto, validazione della copertura CAP, calcolo e selezione dello slot orario disponibile e conferma della prenotazione.

\begin{figure}[H]
    \centering
    % \includegraphics[width=0.85\textwidth]{figure/act_ritiro_domicilio.png}
    \fbox{\parbox[c][6.5cm][c]{0.85\textwidth}{\centering \textsf{[ Inserire qui l'Activity Diagram per la Prenotazione Ritiro a Domicilio esportato da Visual Paradigm ]}}}
    \caption{Activity Diagram --- Processo di Prenotazione Ritiro a Domicilio}
    \label{fig:act_ritiro}
\end{figure}

\subsubsection{Activity Diagram: Esecuzione Ritiro e Registrazione Esito (Autista)}
Il diagramma modella il ciclo operativo dell'autista, dalla presa in carico dell'itinerario giornaliero, al raggiungimento dell'indirizzo e registrazione conclusiva dell'esito.

\begin{figure}[H]
    \centering
    % \includegraphics[width=0.85\textwidth]{figure/act_autista.png}
    \fbox{\parbox[c][6.5cm][c]{0.85\textwidth}{\centering \textsf{[ Inserire qui l'Activity Diagram per l'Esecuzione Ritiro Autista esportato da Visual Paradigm ]}}}
    \caption{Activity Diagram --- Flusso Esecuzione Ritiro e Chiusura Intervento}
    \label{fig:act_autista}
\end{figure}

\subsubsection{Activity Diagram: Accettazione e Conferimento al Varco (Operatore di Sede)}
Il diagramma descrive il flusso di verifica del pass di prenotazione, ispezione visiva dei rifiuti e convalida dello scarico presso l'isola ecologica.

\begin{figure}[H]
    \centering
    % \includegraphics[width=0.85\textwidth]{figure/act_conferimento_sede.png}
    \fbox{\parbox[c][6.5cm][c]{0.85\textwidth}{\centering \textsf{[ Inserire qui l'Activity Diagram per il Conferimento in Sede esportato da Visual Paradigm ]}}}
    \caption{Activity Diagram --- Flusso Accettazione e Verifica al Varco}
    \label{fig:act_sede}
\end{figure}

% =========================================================================
% 5.2 SEQUENCE DIAGRAMS
% =========================================================================
\newpage
\subsection{Sequence Diagrams}
I Sequence Diagram descrivono l'interazione temporale tra gli oggetti del sistema durante l'esecuzione dei casi d'uso, applicando il pattern architetturale \textbf{BCE} (\textit{Boundary}, \textit{Control}, \textit{Entity}).

\subsubsection{Sequence Diagram: UC-C01 Prenotazione Ritiro a Domicilio}
Descrive la collaborazione tra l'interfaccia utente (\textit{Boundary} \texttt{PrenotazioneRitiroUI}), il coordinatore logico (\textit{Control} \texttt{GestorePrenotazioneRitiro}) e le entità di dominio (\textit{Entity} \texttt{Cittadino}, \texttt{ZonaCAP}, \texttt{Veicolo}, \texttt{Prenotazione}).

\begin{figure}[H]
    \centering
    % \includegraphics[width=0.9\textwidth]{figure/seq_prenota_ritiro.png}
    \fbox{\parbox[c][7.5cm][c]{0.85\textwidth}{\centering \textsf{[ Inserire qui il Sequence Diagram (BCE) per la Prenotazione Ritiro esportato da Visual Paradigm ]}}}
    \caption{Sequence Diagram (BCE) --- Realizzazione UC-C01 (Prenotazione Ritiro a Domicilio)}
    \label{fig:seq_ritiro}
\end{figure}

\subsubsection{Sequence Diagram: UC-A02 Registrazione Esito Ritiro}
Descrive la collaborazione tra l'applicazione mobile dell'autista (\textit{Boundary} \texttt{PannelloAutistaUI}), il controllore degli esiti (\textit{Control} \texttt{GestoreEsitoRitiro}) e le entità (\textit{Entity} \texttt{Prenotazione}, \texttt{Itinerario}, \texttt{Notifica}).

\begin{figure}[H]
    \centering
    % \includegraphics[width=0.9\textwidth]{figure/seq_esito_ritiro.png}
    \fbox{\parbox[c][7.5cm][c]{0.85\textwidth}{\centering \textsf{[ Inserire qui il Sequence Diagram (BCE) per la Registrazione Esito Ritiro esportato da Visual Paradigm ]}}}
    \caption{Sequence Diagram (BCE) --- Realizzazione UC-A02 (Registrazione Esito Ritiro)}
    \label{fig:seq_esito}
\end{figure}

\subsubsection{Sequence Diagram: UC-O02 Convalida e Scarico al Varco}
Descrive la scansione del pass di conferimento presso il varco (\textit{Boundary} \texttt{PannelloVarcoUI}), l'elaborazione del controllore (\textit{Control} \texttt{GestoreAccettazioneSede}) e l'aggiornamento dell'entità (\textit{Entity} \texttt{Prenotazione}, \texttt{SedeAMA}).

\begin{figure}[H]
    \centering
    % \includegraphics[width=0.9\textwidth]{figure/seq_conferimento_sede.png}
    \fbox{\parbox[c][7.5cm][c]{0.85\textwidth}{\centering \textsf{[ Inserire qui il Sequence Diagram (BCE) per la Convalida Conferimento esportato da Visual Paradigm ]}}}
    \caption{Sequence Diagram (BCE) --- Realizzazione UC-O02 (Convalida Conferimento al Varco)}
    \label{fig:seq_conferimento}
\end{figure}

% =========================================================================
% 5.3 CLASS DIAGRAMS
% =========================================================================
\newpage
\subsection{Class Diagrams}
I Class Diagram rappresentano la struttura statica delle entità, le classi di controllo e interfaccia, gli attributi, i metodi e le relazioni (associazioni, aggregazioni, composizioni ed ereditarietà).

\subsubsection{Class Diagram Unrefined (Modello Concettuale di Dominio)}
Il modello \textit{Unrefined} illustra le entità concettuali e le loro relazioni primarie derivanti direttamente dall'analisi del dominio applicativo e dal glossario.

\begin{figure}[H]
    \centering
    % \includegraphics[width=0.95\textwidth]{figure/class_unrefined.png}
    \fbox{\parbox[c][8.5cm][c]{0.9\textwidth}{\centering \textsf{[ Inserire qui il Class Diagram Unrefined esportato da Visual Paradigm ]}}}
    \caption{Class Diagram --- Versione Unrefined (Modello di Dominio)}
    \label{fig:class_unrefined}
\end{figure}

\subsubsection{Class Diagram Refined (Modello di Specifica Consolidato)}
Il modello \textit{Refined} integra le classi Boundary e Control emerse dai Sequence Diagram, esplicita le visibilità (\textit{public}, \textit{protected}, \textit{private}), le tipizzazioni dei parametri e dei valori di ritorno, le molteplicità delle associazioni e la navigabilità.

\begin{figure}[H]
    \centering
    % \includegraphics[width=0.95\textwidth]{figure/class_refined.png}
    \fbox{\parbox[c][8.5cm][c]{0.9\textwidth}{\centering \textsf{[ Inserire qui il Class Diagram Refined esportato da Visual Paradigm ]}}}
    \caption{Class Diagram --- Versione Refined (Specifica OOA Consolidata)}
    \label{fig:class_refined}
\end{figure}
